\documentclass[11pt]{article}
\usepackage[utf8]{inputenc}
\usepackage{amsmath,amssymb}
\usepackage{graphicx}
\usepackage{booktabs}
\usepackage{authblk}
\usepackage[margin=1in]{geometry}
\usepackage[hidelinks]{hyperref}
\usepackage{xurl}
\usepackage{setspace}

\title{Elementary Charges as Configurations, not Objects:\\
Why RealQM Needs No Relativity}
\author[1]{Claes Johnson}
\affil[1]{KTH Royal Institute of Technology, SE-100 44 Stockholm, Sweden \\
\texttt{claesjohnson@gmail.com}}
\date{\today \\[0.4em]
\small\textit{Argument by the author; drafting prepared with AI assistance (Claude, Anthropic), reviewed and verified by the author.}}

\begin{document}
\maketitle

\begin{abstract}
We take a single ontological step and follow its consequences: in RealQM an electron or a proton is not a
standalone object but a \emph{domain of a configuration} --- a region of one coupled charge-density
solution, its boundary co-determined by every other charge. On this reading an isolated electron or proton
does not exist; only bound configurations (atoms, molecules, nuclei) do. Three things follow. First, the
apparent defect that a free charge in RealQM merely \emph{spreads out} --- the Coulomb-plus-kinetic
functional has no localised free-particle minimiser --- is not a defect but the theory correctly reporting
that there is no such object; every actual charge gets its size from Coulomb balance \emph{within} a
configuration. Second, RealQM then needs \emph{no relativity}: the only thing that ever demanded a
rest-energy or Compton scale was giving an \emph{isolated} particle an intrinsic size, and with isolated
particles removed the demand vanishes --- all scales come from Coulomb configurations, mass is at most a
scale-label, and $E=mc^2$ is never invoked. Third, the electron's magnetic moment follows without
relativity as a \emph{configuration property}: a weakly-confined electron (atom, trap) can afford the
quantised current-winding that gives $\mu_B$, while a tightly-caged one (nucleus) cannot and carries none,
so nuclear moments sit at nuclear-magneton scale. The one residue this leaves is not relativity but a
sharp structural question --- why the winding's gyromagnetic factor is $2$ rather than $1$. The picture is
consistent with the rest of the programme (real 3D densities, size-not-mass, the neutron as $p+e$), and it
closes the free-charge problem by dissolving the free charge.
\end{abstract}

\setstretch{1.12}

\section{The step: a charge is a configuration, not an object}
\label{sec:step}

In RealQM~\cite{RealQM} the electrons and protons of a system are one-electron/one-proton charge densities
$\psi_i$ supported on non-overlapping domains $\Omega_i\subset\mathbb R^3$, the boundaries $\Gamma_i$
determined \emph{jointly} with the $\psi_i$ by minimisation of the Coulomb energy. An individual
$\psi_i$ is therefore not an independent thing: it is a \emph{region of one coupled solution}, its extent
and shape fixed by the presence of all the other charges. We take this literally. An electron or a proton
\emph{is} such a domain-of-a-configuration; it has no existence apart from the configuration it partakes
in. The fundamental existents are atoms, molecules, and nuclei --- bound configurations of charge --- and
``the electron'' and ``the proton'' are abstractions read off them, not standalone objects that pre-exist
and are later assembled.

This is not a new postulate grafted onto RealQM; it is what the formalism already implies, made explicit.
It is also the natural end of positions already taken in the programme: that the neutron \emph{is} a bound
$p+e$ pair rather than an elementary particle, and that a bound nucleon has no separate identity. Here we
apply the same reading to the free case, where it has its sharpest consequence.

\section{The free charge dissolves --- and with it, a false problem}
\label{sec:free}

A lone electron in RealQM, with no other charge to attract it, minimises $\tfrac12\int|\nabla\psi|^2$
alone; the infimum is zero, approached by spreading $\psi$ over ever larger volume, and there is
\emph{no localised minimiser}. The same holds for a lone proton. Read as a statement about objects, this
looks like a failure: the theory ``cannot confine'' a free particle. Read through
\S\ref{sec:step}, it is not a failure but a \emph{correct report}: there is no isolated electron for the
theory to confine. The spreading is the theory saying, of a would-be free charge, \emph{no such object}.

Every charge we actually deal with is in a configuration, and takes its size from the Coulomb balance
there: the Bohr scale for an electron bound to a nucleus, the femtometre scale for a proton or electron in
a nucleus, the trap scale for an electron held in a Penning trap. Even the experimentalist's ``free''
electron is an electron weakly coupled to the fields of an apparatus --- a configuration, not a bare object
adrift in empty space. A genuinely field-free electron in vacuum would disperse; we never have one. The
size of a charge is thus always a property of a configuration, never of an isolated particle.

\section{No isolated particle, no need for relativity}
\label{sec:norel}

The consequence for relativity is direct. The one place a \emph{rest-energy} or \emph{Compton} scale was
forced upon us was in giving an \emph{isolated} particle an intrinsic size: a free electron, if it is a
real object, needs \emph{something} to fix its extent, and standard physics supplies the Compton
wavelength $\hbar/m_e c$ set by the rest mass $m_e c^2$ --- a relativistic quantity. Remove the isolated
particle and the demand disappears. Sizes come from Coulomb configurations; mass enters, if at all, only
as a \emph{scale-label} selecting which regime (atomic or nuclear) a configuration sits in, exactly the
``size, not mass'' reading the programme already adopts~\cite{RealNucleus}. There is no free-particle
rest energy to account for, so $E=mc^2$ is never invoked, and RealQM stays what it is: a non-relativistic
theory of Coulomb charge densities in real space. Relativity, on this view, is not \emph{denied} --- it is
simply \emph{never needed}, because the objects whose intrinsic properties required it do not exist.

\section{The magnetic moment as a configuration property}
\label{sec:moment}

The magnetic moment makes the point concretely, and without relativity. In the current-bearing
(complex, time-dependent) form of RealQM~\cite{ComplexRealQM}, the moment is that of the charge current,
$\boldsymbol\mu = \tfrac12\int\mathbf r\times\mathbf J$, and a phase that \emph{winds in space} carries a
standing current with, for a normalised electron density, the quantised value $\mu_z = -m\,\mu_B$ ---
\emph{independent of the density's size}. Whether a given configuration \emph{can} support such a winding
is an energetic question: the winding costs $\Delta E \sim \hbar^2/(m_e\sigma^2)$, cheap for a large
confinement scale $\sigma$ and prohibitive for a small one~\cite{NuclearMoment}. Hence:
\begin{itemize}
\item a \emph{weakly-confined} electron --- in an atom, or in a trap, at large $\sigma$ --- can afford the
winding, and carries $\mu_z = -\mu_B$; this is the ``free electron's moment,'' now a property of its
(loose) configuration, not of an isolated object;
\item a \emph{tightly-caged} electron --- in a nucleus, at $\sigma\sim$~fm --- cannot afford it (the cost
runs to GeV, thousands of times the binding), settles in the winding-free state, and carries \emph{no}
moment, so nuclear moments come out at nuclear-magneton scale.
\end{itemize}
The Bohr magneton is thus the quantised winding a loose configuration permits, not an intrinsic
rest-energy attribute; the Penning-trap electron carries it for the same reason an atomic electron can.
No Compton scale, no relativity, enters anywhere.

\paragraph{Radiation, likewise, is charge moving in space.} The companion point about radiation is of the
same character. Radiation is not a ``superposition of stationary states'' --- that is configuration-space
language --- but the \emph{physical oscillation of a charge density in real space}: a charge cloud that
sloshes is an antenna and radiates by the Larmor formula, driven by the current $\mathbf J$. Oscillation
gives radiation and circulation gives a magnetic moment; both are charge in motion in three-dimensional
space, within a configuration, and neither needs relativity or an isolated particle to carry it.

\section{The residue: a factor of two, not relativity}
\label{sec:residue}

This dissolves three things at once --- the existence of the free charge, the need for its intrinsic size,
and the need for relativity. It does not, by itself, deliver the electron's \emph{spin} gyromagnetic
factor. An orbital winding has $g=1$ ($\mu = \mu_B$ for one unit of angular momentum), whereas the
electron's measured moment corresponds to $g=2$, with a further small anomaly $g/2 = 1.00116\ldots$. So a
residue remains --- but note its character has changed. It is no longer the intractable
\emph{``how does a free particle acquire an intrinsic relativistic moment?''}, unanswerable within a
non-relativistic frame; it is the sharp, bounded structural question \emph{``why is the gyromagnetic
factor of the winding current $2$ and not $1$?''} --- a question about the pattern of the current in a
configuration, to be settled inside RealQM, not by importing relativity. We flag it as the open problem,
and claim only that it is now the \emph{right size} of problem.

\section{Consistency and conclusion}
\label{sec:concl}

The reading coheres with the rest of the programme. It is the same realism as $\psi^2$-as-charge-density
and the Madelung flow made physical in real space; it is the same ``size, not mass''; it is the same
denial of independent identity to a bound nucleon, now extended to the free case. What it buys is a clean
closure of the free-charge problem: RealQM cannot confine a free electron or proton because there is no
such object to confine, and precisely for that reason it needs no rest-energy scale and no relativity ---
every size and every moment is a property of a Coulomb configuration in three-dimensional space. Is it
possible to do physics this way? For everything the programme has treated --- binding, geometry, reactions,
nuclear structure, the caged electron's vanishing moment, radiation as moving charge --- it appears not
only possible but natural. What remains is a factor of two in the electron's spin moment, which is a
question about currents, not about relativity.

\begin{thebibliography}{9}
\bibitem{RealQM} C.~Johnson, \emph{Real Quantum Mechanics as Multiphase 3D Continuum Mechanics},
manuscript, 2026; \url{https://claes542.github.io/RealMolecule/gallery.html}.
\bibitem{RealNucleus} C.~Johnson, \emph{RealNucleus: Nuclear Binding as Dual Coulomb Confinement, without
a Strong or Weak Force}, 2026.
\bibitem{ComplexRealQM} C.~Johnson, \emph{Complex Time-Dependent RealQM: Currents, Radiation, and
Magnetism from the Real-Space Continuum}, 2026.
\bibitem{NuclearMoment} C.~Johnson, \emph{Magnetism as Circulating Charge: Why the Nuclear Electron
Carries No Magnetic Moment}, 2026.
\end{thebibliography}

\end{document}
